\chapter{Hypervisors, Virtual Machines and Containers}\label{ch:virtualisation}
\bp{1.2}

\begin{outcomes}
By the end of this chapter you will be able to:
\begin{tight}
  \item \textbf{Distinguish} a type~1 from a type~2 hypervisor and say where each
        belongs.
  \item \textbf{Describe} what a virtual machine contains and how a virtual
        switch connects it to the physical network.
  \item \textbf{Compare} virtual machines with containers on isolation, start
        time, size and what they share.
  \item \textbf{Configure} the physical switch port a hypervisor host attaches to
        --- which is the part of this topic that is actually your job.
\end{tight}
\end{outcomes}

\begin{prereq}
\chref{vlans} and \chref{trunking} --- a hypervisor uplink is a trunk, and this
chapter will not make sense without knowing what one is. \chref{l2security} for
port security, which is the wrong tool on such a port. \chref{ethernet} for MAC
learning.
\end{prereq}

\begin{keyterms}
hypervisor $\bullet$ type 1 $\bullet$ type 2 $\bullet$ virtual machine $\bullet$
guest OS $\bullet$ vNIC $\bullet$ virtual switch $\bullet$ port group $\bullet$
container $\bullet$ image $\bullet$ kernel $\bullet$ orchestration $\bullet$
east-west traffic $\bullet$ live migration
\end{keyterms}

\begin{examnote}[\emph{Describe the role and function}]
Topic~\tp{1.2} is \emph{describe the role and function of hypervisors, virtual
machines, and containers}. No configuration of any hypervisor is required, and no
vendor product is named.

But it sits in domain~1, \emph{Network Infrastructure and Connectivity}, and not
by accident. A CCNA is not asked to administer a hypervisor; a CCNA is asked to
\textbf{connect one correctly}, and to recognise the several things a virtualised
server does to a network that a physical server never did. Section~\ref{sec:vnetimpact} is where
those live, and it is the part of this chapter you will use at work next week.
\end{examnote}

\section{Why any of this exists}

A physical server running one application typically uses 5 to 15\% of its
capacity. The other 85\% is bought, powered, cooled and wasted. Virtualisation
puts many logical servers on one physical one, each believing it has the machine
to itself.

\begin{definitionbox}[Hypervisor]
Software that presents virtual hardware --- CPU, memory, disk, network interface
--- to several guest operating systems at once, and schedules the real hardware
between them. Also called a \term{virtual machine monitor}.
\end{definitionbox}

\subsection{Type 1 and type 2}

\ccnafigh{%
\begin{tikzpicture}[font=\scriptsize]
  \tikzset{
    lay/.style={draw=#1, fill=#1!10, text=#1!40!black, minimum height=0.62cm,
                align=center, font=\tiny\bfseries, line width=0.7pt},
    ttl/.style={font=\scriptsize\bfseries, text=neutral!45!black}}

  % Type 1
  \node[ttl] at (3.3,3.35) {Type 1 --- bare metal};
  \node[lay=neutral, minimum width=6.6cm] at (3.3,0.0) {Physical hardware};
  \node[lay=dom3,    minimum width=6.6cm] at (3.3,0.75) {Hypervisor (ESXi, Hyper-V, KVM, Xen)};
  \foreach \x/\n in {1.25/VM, 3.3/VM, 5.35/VM} {
    \node[lay=dom1, minimum width=1.9cm] at (\x,1.6) {Guest OS};
    \node[lay=dom1, fill=dom1!22, minimum width=1.9cm] at (\x,2.3) {App};
  }
  \node[font=\tiny, text=neutral!85!black, align=center] at (3.3,2.85)
      {data centre, production};

  % Type 2
  \node[ttl] at (11.0,3.35) {Type 2 --- hosted};
  \node[lay=neutral, minimum width=5.4cm] at (11.0,0.0)  {Physical hardware};
  \node[lay=dom2,    minimum width=5.4cm] at (11.0,0.75) {Host OS (Windows, macOS, Linux)};
  \node[lay=dom3,    minimum width=5.4cm] at (11.0,1.5)  {Hypervisor (Workstation, VirtualBox)};
  \foreach \x in {9.7, 12.3} {
    \node[lay=dom1, minimum width=2.2cm] at (\x,2.3) {Guest OS + App};
  }
  \node[font=\tiny, text=neutral!85!black, align=center] at (11.0,2.85)
      {laptop, lab, your CML};
\end{tikzpicture}}%
{The difference is one layer. A type~1 hypervisor \emph{is} the operating system;
a type~2 hypervisor is an application running on one.}{hypervisors}

\begin{longtable}{@{}L{3.0cm}L{5.8cm}L{6.1cm}@{}}
\toprule
\rowcolor{primary!15}
\textbf{} & \textbf{Type 1 (bare metal, native)} & \textbf{Type 2 (hosted)} \\
\midrule
\endhead
Runs on & The hardware directly &
On top of a normal host operating system \\
\rowcolor{lightbg}
Performance & Higher --- no host OS in the path &
Lower --- every operation passes through the host OS \\
Examples & VMware ESXi, Microsoft Hyper-V, KVM, Xen &
VMware Workstation and Fusion, Oracle VirtualBox, Parallels \\
\rowcolor{lightbg}
Used for & Production servers, data centres, cloud &
Development, testing, labs, running one OS inside another \\
Boots into & The hypervisor itself &
The host OS; the hypervisor is launched afterwards \\
\bottomrule
\end{longtable}

\begin{conceptbox}[You have almost certainly used both]
The CML or GNS3 you run the labs in this book on is a type~2 hypervisor, or runs
inside one. The cloud instances your university's systems run on are type~1.

If an exam question describes a hypervisor installed \emph{onto} a server with no
other operating system, it is type~1. If it describes one installed \emph{into}
Windows or macOS, it is type~2. That is the whole distinction.
\end{conceptbox}

\section{What a virtual machine is}

A VM is a set of files --- a configuration file describing its virtual hardware,
and one or more virtual disk files --- plus the state the hypervisor keeps while
it runs. It contains a complete guest operating system with its own kernel,
drivers, and boot process.

Because it is files, it can be copied, snapshotted, moved between hosts while
running (\term{live migration}, or vMotion), and restored. Those conveniences are
also what makes it a networking problem, which is Section~\ref{sec:vnetimpact}.

\subsection{The virtual switch}

Every VM has one or more \term{vNICs}, each with its own MAC address. Those
vNICs plug into a \term{virtual switch} inside the hypervisor, and the virtual
switch has one or more physical NICs as uplinks to the real network.

\ccnafigh{%
\begin{tikzpicture}[font=\scriptsize]
  \tikzset{
    vm/.style={rounded corners=2pt, draw=dom1, fill=dom1!8, text=dom1!45!black,
               line width=0.8pt, minimum width=1.75cm, minimum height=0.95cm,
               align=center, font=\tiny\bfseries},
    w/.style={line width=0.85pt, neutral!75}}

  % Vertical budget, so nothing collides: host box -0.55 to 3.05;
  % its own label occupies the top-right corner, the east-west arc and
  % its caption the band from 1.9 to 2.7, the VMs 0.88 to 1.83, and the
  % virtual switch -0.38 to 0.28.
  \draw[dom3!45, fill=dom3!4, line width=1pt, rounded corners=4pt]
      (-0.5,-0.55) rectangle (8.6,3.05);
  \node[font=\tiny\bfseries, text=dom3!45!black, anchor=north east]
      at (8.5,2.98) {HYPERVISOR HOST};

  \node[vm] (v1) at (0.9,1.35) {VM 1\\\tiny\mdseries VLAN 10};
  \node[vm] (v2) at (3.2,1.35) {VM 2\\\tiny\mdseries VLAN 10};
  \node[vm] (v3) at (5.5,1.35) {VM 3\\\tiny\mdseries VLAN 20};
  \node[vm] (v4) at (7.7,1.35) {VM 4\\\tiny\mdseries VLAN 30};

  \node[rounded corners=2pt, draw=dom2, fill=dom2!10, text=dom2!45!black,
        line width=0.9pt, minimum width=8.0cm, minimum height=0.66cm,
        font=\tiny\bfseries] (vs) at (4.0,-0.05) {VIRTUAL SWITCH \quad
        \mdseries port groups: VLAN 10 $\cdot$ VLAN 20 $\cdot$ VLAN 30};

  \foreach \v in {v1,v2,v3,v4} \draw[w] (\v) -- (\v |- vs.north);

  \node[rounded corners=3pt, draw=dom2, fill=dom2!8, text=dom2!45!black,
        line width=0.9pt, minimum width=1.5cm, minimum height=1.0cm,
        align=center, font=\scriptsize\bfseries] (sw) at (11.9,-0.05)
        {\faIcon{sitemap}\\SW1};

  % pos=0.5 keeps the two-line label clear of the host box's right edge.
  \draw[w, dash pattern=on 3pt off 2pt] (vs.east) -- (sw.west)
      node[font=\tiny, text=dom2!45!black, pos=0.5, above=2pt, align=center]
      {\textbf{802.1Q trunk}\\VLANs 10, 20, 30};

  % East--west traffic arcs over the two VMs rather than through them.
  \draw[{Stealth[length=2mm]}-{Stealth[length=2mm]}, dom4!70, line width=0.9pt]
      (v1.north) to[bend left=45] (v2.north);
  \node[font=\tiny, text=dom4!45!black, align=center, anchor=south]
      at (2.05,2.28) {east--west: never leaves the host};
\end{tikzpicture}}%
{VM 1 talking to VM 2 never reaches SW1. Everything else does, over a trunk --- a
single physical port carrying three VLANs.}{vswitch}

\begin{definitionbox}[Port group]
A named policy on the virtual switch --- a VLAN ID, security settings and traffic
shaping --- that VMs are attached to. It is the virtual equivalent of deciding
which VLAN a switch port belongs to, and the VLAN ID configured on it must match
one the physical switch actually trunks.
\end{definitionbox}

\section{Containers}

A container does not virtualise the hardware. It virtualises the \emph{operating
system}: several isolated user-space environments share one kernel.

\ccnafigh{%
\begin{tikzpicture}[font=\scriptsize]
  \tikzset{
    lay/.style={draw=#1, fill=#1!10, text=#1!40!black, minimum height=0.62cm,
                align=center, font=\tiny\bfseries, line width=0.7pt},
    ttl/.style={font=\scriptsize\bfseries, text=neutral!45!black}}

  \node[ttl] at (3.3,3.35) {Virtual machines};
  \node[lay=neutral, minimum width=6.6cm] at (3.3,0.0)  {Physical hardware};
  \node[lay=dom3,    minimum width=6.6cm] at (3.3,0.72) {Hypervisor};
  \foreach \x in {1.25, 3.3, 5.35} {
    \node[lay=dom1, minimum width=1.9cm] at (\x,1.44) {\textbf{Guest OS}};
    \node[lay=dom1, fill=dom1!22, minimum width=1.9cm] at (\x,2.16) {App + libs};
  }
  \node[font=\tiny, text=accent!70!black, align=center] at (3.3,2.82)
      {GB in size, seconds to minutes to boot};

  \node[ttl] at (11.0,3.35) {Containers};
  \node[lay=neutral, minimum width=6.6cm] at (11.0,0.0)  {Physical hardware};
  \node[lay=dom2,    minimum width=6.6cm] at (11.0,0.72) {\textbf{Host OS --- one shared kernel}};
  \node[lay=dom3,    minimum width=6.6cm] at (11.0,1.44) {Container engine};
  \foreach \x in {8.95, 11.0, 13.05} {
    \node[lay=dom1, fill=dom1!22, minimum width=1.9cm] at (\x,2.16) {App + libs};
  }
  \node[font=\tiny, text=dom2!45!black, align=center] at (11.0,2.82)
      {MB in size, milliseconds to start};
\end{tikzpicture}}%
{One layer removed. Containers share the host's kernel, which is why they are
small and fast --- and why the isolation between them is weaker.}{containers}

\begin{longtable}{@{}L{3.0cm}L{5.8cm}L{6.1cm}@{}}
\toprule
\rowcolor{primary!15}
\textbf{} & \textbf{Virtual machine} & \textbf{Container} \\
\midrule
\endhead
Contains & A full guest OS with its own kernel &
An application and its libraries only \\
\rowcolor{lightbg}
Shares & Nothing above the hypervisor &
\textbf{The host's kernel} \\
Size & Gigabytes & Megabytes \\
\rowcolor{lightbg}
Start time & Seconds to minutes --- it boots &
Milliseconds --- there is nothing to boot \\
Isolation & Strong. A compromised guest is still inside its VM &
Weaker. A kernel vulnerability is shared by every container on the host \\
\rowcolor{lightbg}
Guest OS choice & Any, including a different family from the host &
Must be compatible with the host kernel \\
Typical use & Consolidating servers; running different operating systems &
Packaging and scaling one application; microservices \\
\rowcolor{lightbg}
Managed by & The hypervisor, and tools such as vCenter &
A container engine, and an orchestrator such as Kubernetes \\
\bottomrule
\end{longtable}

\begin{conceptbox}[They are not competitors]
The common real deployment is containers \emph{inside} virtual machines: the VM
provides hard isolation and the ability to move workloads between hosts, and the
container provides fast, repeatable application packaging inside it. Every major
cloud provider runs some version of this.

An exam question asking you to choose between them is asking about the properties
in the table --- kernel sharing, size, start time, isolation --- not about which
is better.
\end{conceptbox}

\begin{definitionbox}[Orchestration]
Running one container by hand is easy. Running four hundred across thirty hosts,
restarting the ones that fail, scaling with demand and routing traffic to
whichever are healthy is not. An \term{orchestrator} --- Kubernetes is the
dominant one --- does that scheduling, health-checking and service discovery.

For this exam, know what the word means and what problem it solves.
\end{definitionbox}

\section{What this does to your network}\label{sec:vnetimpact}

This is the section that matters to a network engineer, and it is where the
blueprint's placement of this topic in domain~1 makes sense.

\begin{longtable}{@{}L{4.6cm}L{10.4cm}@{}}
\toprule
\rowcolor{primary!15}
\textbf{Consequence} & \textbf{What it means for you} \\
\midrule
\endhead
\textbf{The uplink is a trunk} &
One physical port carries every VLAN the host's VMs live in. Configuring it as an
access port strands every VM outside that one VLAN \\
\rowcolor{lightbg}
\textbf{Many MACs on one port} &
Dozens, changing as VMs start and stop. Port security on that port is simply
the wrong tool \\
\textbf{MACs move without anyone touching a cable} &
Live migration relocates a running VM to another host, and its MAC appears on a
different switch port seconds later. This looks exactly like a loop \\
\rowcolor{lightbg}
\textbf{East--west traffic is invisible} &
Two VMs on the same host and VLAN never reach the physical switch. You cannot
see it, filter it with a physical ACL, or capture it from the network \\
\textbf{The server team owns a switch now} &
The virtual switch's VLAN configuration is set by whoever administers the
hypervisor. A mismatch between their port group and your trunk is a shared fault
with two owners \\
\rowcolor{lightbg}
\textbf{Nothing has a fixed location} &
Documentation naming a physical port for a service goes stale immediately \\
\bottomrule
\end{longtable}

\begin{config}{the physical switch port a hypervisor host attaches to}
interface GigabitEthernet1/0/24
 description ESXi host 3 -- vmnic0
 switchport trunk encapsulation dot1q
 switchport mode trunk
 switchport trunk allowed vlan 10,20,30
 switchport nonegotiate
!
! No port security: a hypervisor presents many MAC addresses and they
! change as machines start, stop and migrate.
!
! PortFast trunk, because the far side is a virtual switch that does
! not run spanning tree and will never send a BPDU.
 spanning-tree portfast trunk
!
! BPDU guard is still correct: nothing on this port should ever send
! a BPDU, so one arriving means something is wrong.
 spanning-tree bpduguard enable
\end{config}

\begin{alertbox}[The virtual switch does not run spanning tree]
A vSwitch is not a learning switch in the ordinary sense: it does not forward
traffic between its uplinks, so it cannot create a loop, and it does not
participate in STP.

Two consequences. \cmd{spanning-tree portfast trunk} is appropriate on the
physical port, because there is no risk of a loop and waiting 30 seconds on every
host reboot is pure cost. And \textbf{BPDU guard should still be enabled} ---
precisely because a correctly configured vSwitch will never send a BPDU, one
arriving means somebody has attached something else.
\end{alertbox}

\begin{pitfall}
\begin{tight}
  \item \textbf{Access port on a hypervisor uplink.} Everything in one VLAN
        works, everything else is stranded, and the fault looks like a VM
        problem.
  \item \textbf{Port security on a hypervisor uplink.} Err-disables the port and
        takes down every VM on the host.
  \item \textbf{Reading migration MAC moves as a loop.} A MAC flap between two
        \emph{access-layer uplinks} at the moment a VM migrated is expected
        behaviour.
  \item \textbf{Forgetting to add a VLAN to \cmd{allowed}} when the server team
        adds a port group. Their side is right, yours is not, and neither team
        sees both halves.
  \item \textbf{Expecting to capture east--west traffic} on the physical switch.
        It never gets there.
  \item \textbf{Assuming containers are just small VMs.} The shared kernel is the
        whole difference, and it is where the exam question will be.
  \item \textbf{Confusing type 1 with type 2.} Type 1 replaces the OS; type 2
        runs on one.
\end{tight}
\end{pitfall}

\begin{breakfix}[A new hypervisor host is racked. VMs in VLAN 10 work; VMs in VLANs 20 and 30 are dead.]
The server team reports that of twelve VMs on the new host, the four in VLAN~10
have full connectivity and the eight in VLANs 20 and 30 cannot reach anything,
including their own default gateways. The vSwitch configuration is identical to
the three existing hosts, which all work. Later in the morning the port
err-disables entirely and all twelve VMs drop.

\begin{verify}{SW1}{show interfaces GigabitEthernet1/0/24 switchport}
Name: Gi1/0/24
Switchport: Enabled
Administrative Mode: static access
Operational Mode: static access
Access Mode VLAN: 10 (STAFF)
Trunking Native Mode VLAN: 1 (default)
\end{verify}

\begin{verify}{SW1}{show interfaces status err-disabled}
Port      Name               Status       Reason               Err-disabled Vlans
Gi1/0/24  ESXi host 4        err-disabled psecure-violation
\end{verify}

\textbf{Diagnosis.} Two faults on one port, and the first explains the pattern
while the second explains the timing.

\textbf{The port is an access port in VLAN 10.} A hypervisor uplink must be a
trunk. VMs in the port group tagged VLAN~10 happen to work because the vSwitch
sends their frames untagged and the access port puts them in VLAN~10 --- the
right answer by accident. Frames from VLANs 20 and 30 arrive tagged, and an
access port discards tagged frames for other VLANs. Eight VMs are stranded.

\textbf{The port also has port security.} A hypervisor presents one MAC per vNIC
plus the host's own management address --- thirteen or more on this host. As the
remaining VMs powered on through the morning, the maximum was exceeded and the
port err-disabled, which is why the working four then failed too.

The tell that this is a switch-port fault and not a hypervisor fault is that the
vSwitch configuration matches three working hosts. When one member of an
identical set behaves differently, look at what is \emph{not} identical --- which
here is the physical port it was patched into.

\textbf{Fix.}

\begin{config}{correcting the uplink}
interface GigabitEthernet1/0/24
 description ESXi host 4 -- vmnic0
 no switchport port-security
 no switchport access vlan 10
 switchport trunk encapsulation dot1q
 switchport mode trunk
 switchport trunk allowed vlan 10,20,30
 switchport nonegotiate
 spanning-tree portfast trunk
 spanning-tree bpduguard enable
!
 shutdown
 no shutdown
\end{config}

The \cmd{shutdown} then \cmd{no shutdown} is required: removing the port security
configuration does not clear an existing \cmd{err-disabled} state.

\textbf{The process lesson.} This port was almost certainly configured from the
site's standard access-port template, which is correct for a desk and wrong for a
server. Sites that hit this repeatedly keep a separate documented template for
hypervisor uplinks --- trunk, explicit allowed list, no port security,
\cmd{portfast trunk}, BPDU guard --- and name it in the change request.
\end{breakfix}

\begin{lab}{Connecting a hypervisor properly}{VirtualBox, VMware Workstation, or CML}
\textbf{Platform note.} You do not need a data centre. Any type~2 hypervisor on
your own machine demonstrates every concept here; the switch-side tasks can be
done in Packet Tracer or CML against a simulated host.

\begin{tasks}
  \item Identify the hypervisor you already use for the labs in this book and
        classify it as type~1 or type~2. Justify the classification in one
        sentence.
  \item Create two VMs. Record the MAC address of each vNIC, and confirm they
        differ from the host's physical NIC.
  \item Put both VMs on a \emph{host-only} or internal network. Confirm they can
        reach each other and that this traffic never appears on your physical
        network --- capture on the physical interface to prove it. This is
        east--west traffic.
  \item Switch both to bridged networking. Confirm they now appear on the
        physical network with their own MAC addresses, and find them in your
        switch's MAC address table.
  \item \textbf{Switch side.} Configure the port your host connects to as an
        access port in one VLAN. Put one VM in that VLAN and one in another.
        Reproduce the Break \& Fix symptom, then convert the port to a trunk.
  \item Apply \cmd{switchport port-security maximum 1} to that trunk. Power on
        the second VM and observe. Recover the port.
  \item \textbf{Containers.} Install a container engine and run two containers.
        Compare their start time and image size against your VMs' boot time and
        disk size. Record actual numbers.
  \item Inspect the container network. Find the bridge, the addresses assigned,
        and confirm that outbound traffic is being NATted by the host. Relate
        this to \chref{nat}.
  \item Run \cmd{uname -r} inside two containers and on the host. Explain the
        result, and state what it implies for isolation.
  \item \textbf{Extension.} Write the switch-port template your site should use
        for hypervisor uplinks, with a one-line justification per command, and
        say what is different from your standard access-port template.
\end{tasks}
\end{lab}

\begin{checkpoint}
\begin{tightnum}
  \item What is the single structural difference between a type~1 and a type~2
        hypervisor?
  \item What does a container share with its host that a virtual machine does
        not, and what are the two main consequences?
  \item Why must a hypervisor uplink be a trunk?
  \item Why is port security inappropriate on that port?
  \item What is east--west traffic, and why can you not capture it on the
        physical switch?
  \item A MAC address moves between two switch ports with no cable change. Give a
        legitimate explanation.
\end{tightnum}
\end{checkpoint}

\begin{chaptersummary}
\begin{tight}
  \item A hypervisor presents virtual hardware to several guest operating
        systems. Type~1 runs on the bare metal; type~2 runs on a host OS.
  \item A VM is files: a configuration and virtual disks, containing a complete
        guest OS with its own kernel. It can be copied, snapshotted and migrated
        while running.
  \item Each vNIC has its own MAC and attaches to a virtual switch, whose port
        groups carry VLAN IDs that must match the physical trunk.
  \item Containers share the host kernel: megabytes rather than gigabytes,
        milliseconds rather than seconds, weaker isolation. Orchestrators such as
        Kubernetes schedule them at scale.
  \item The two are commonly combined --- containers inside VMs.
  \item The physical port to a hypervisor is a \textbf{trunk} with an explicit
        allowed list, \textbf{no port security}, \cmd{portfast trunk} and BPDU
        guard.
  \item Live migration moves MAC addresses between ports legitimately. East--west
        traffic never reaches the physical network at all.
\end{tight}
\end{chaptersummary}

\begin{reviewq}
  \item \textbf{(MCQ)} Which describes a type~1 hypervisor?
        \begin{tight}
          \item[(a)] Runs as an application on Windows or macOS
          \item[(b)] Runs directly on the physical hardware with no host OS
          \item[(c)] Shares the host operating system's kernel with its guests
          \item[(d)] Requires an orchestrator such as Kubernetes
        \end{tight}
  \item \textbf{(MCQ)} What do containers share that virtual machines do not?
        \begin{tight}
          \item[(a)] The physical NIC \quad (b)~The host's kernel
          \item[(c)] The hypervisor \quad (d)~A single MAC address
        \end{tight}
  \item \textbf{(Short)} Explain why applying port security to a hypervisor
        uplink causes an outage, and what you should do instead.
  \item \textbf{(Short)} A monitoring team asks why they cannot see traffic
        between two application servers that are known to communicate constantly.
        Give the explanation and one thing that would let them see it.
  \item \textbf{(Applied)} A hypervisor host has VMs in VLANs 10, 20 and 40.
        Write the complete configuration for the switch port it attaches to,
        justify each line, and state which single line would strand two-thirds of
        the VMs if omitted.
  \item \textbf{(Applied)} Compare virtual machines and containers across
        isolation, start time, size and guest OS flexibility, and recommend one
        for each of: running a legacy Windows application, and scaling a web
        front end from 3 to 300 instances during registration week.
\end{reviewq}
